% chapters/ch09_rtos_realtime.tex — Chapter 9: RTOS & Real-Time
% Scope (PLAN.md): task states & scheduling (preemptive/round-robin/
% rate-monotonic), context switch, mutex vs semaphore vs queue, priority
% inversion & inheritance, deadlock, hard vs soft real-time, determinism,
% bare-metal superloop vs RTOS, shared-memory/DPRAM handshake (her BEL work).
% Budget 40 pp, mix 24 Q&A / 12 code / 4 concept.
% All snippets verified gcc -Wall -Wextra -std=c11 (zero warnings).

\chapter{RTOS \& Real-Time}
\label{ch:rtos-realtime}

\begin{partnote}
\textbf{Part C opens here.} \textbf{Scope} --- task states \& scheduling
(preemptive / round-robin / rate-monotonic / EDF) $\rightarrow$ the context
switch $\rightarrow$ mutex vs semaphore vs queue $\rightarrow$ priority
inversion \& inheritance $\rightarrow$ deadlock $\rightarrow$ hard vs soft
real-time \& determinism $\rightarrow$ bare-metal superloop vs RTOS
$\rightarrow$ shared-memory / DPRAM handshake patterns. The DPRAM handshake is
straight off the resume (the BEL real-time DPRAM polling architecture with
sub-1\,ms latency), and priority inversion is \emph{the} classic RTOS
interview question.
\end{partnote}

``Real-time'' does not mean fast --- it means \emph{predictable}: the right
answer by the deadline, every time, even in the worst case. An RTOS exists to
make multi-task timing analysable; the questions here probe whether you
understand the mechanisms (scheduling, synchronisation) \emph{and} their failure
modes (inversion, deadlock, jitter).

%=====================================================================
\section{Primer}
%=====================================================================

\subsection{Tasks, states, and the context switch}

An RTOS runs multiple \textbf{tasks}, each with its own stack and a
\textbf{control block (TCB)} holding its saved registers, stack pointer,
priority, and state. A task is \textbf{running}, \textbf{ready} (runnable, but a
higher-priority task has the CPU), or \textbf{blocked} (waiting on a
semaphore/queue/delay). The scheduler switches the CPU between them with a
\textbf{context switch}: save the running task's registers to its TCB/stack,
load the next task's, and resume --- on Cortex-M this is done in the PendSV
handler, leveraging the hardware exception stacking from Chapter 3. The switch
costs cycles, so its frequency and cost factor into the timing budget.

\subsection{Scheduling policies}

\textbf{Preemptive priority} is the norm in an RTOS: the highest-priority ready
task always runs, preempting lower ones immediately --- this is what makes
response time analysable. \textbf{Round-robin} time-slices equal-priority tasks
for fairness. \textbf{Rate-monotonic scheduling (RMS)} is the classic
fixed-priority rule for periodic tasks: \emph{shorter period $\rightarrow$
higher priority}, which is provably optimal among fixed-priority schemes, with
the Liu--Layland utilisation bound $U \le n(2^{1/n}-1)$ (~69\% as $n\to\infty$)
guaranteeing schedulability. \textbf{Earliest-deadline-first (EDF)} is a dynamic
policy that always runs the task with the nearest deadline and can reach 100\%
utilisation, at the cost of dynamic priorities and harder analysis. Most
embedded RTOSes use fixed-priority preemption with rate-monotonic assignment.

\subsection{Synchronisation: mutex, semaphore, queue}

Three primitives, three jobs. A \textbf{mutex} provides \emph{mutual exclusion}
--- it has an owner, only the owner unlocks it, and it protects a shared
resource (a peripheral, a data structure). A \textbf{semaphore} is a counter for
\emph{signalling/resource-counting} --- a binary semaphore signals an event
(ISR$\rightarrow$task), a counting semaphore guards $N$ identical resources;
it has no owner. A \textbf{queue} passes \emph{data} between tasks (or
ISR$\rightarrow$task) safely, combining buffering and synchronisation --- the
preferred way to move data because it avoids shared mutable state entirely. The
rule of thumb: protect a resource $\rightarrow$ mutex; signal an event
$\rightarrow$ semaphore; move data $\rightarrow$ queue.

\subsection{Priority inversion, deadlock, and DPRAM}

\textbf{Priority inversion}: a high-priority task blocks on a mutex held by a
low-priority task, and a \emph{medium} task (needing no mutex) preempts the low
task --- so the medium task effectively delays the high task indefinitely. The
fix is \textbf{priority inheritance}: temporarily raise the mutex holder to the
waiter's priority so it runs, releases, and unblocks the high task (this is the
Mars Pathfinder bug). \textbf{Deadlock} needs four simultaneous conditions
(mutual exclusion, hold-and-wait, no preemption, circular wait); break any one
--- commonly by always acquiring multiple locks in a \emph{fixed global order}
--- and deadlock is impossible. \textbf{DPRAM} (dual-port RAM) lets two
processors share memory; a \textbf{flag/mailbox handshake} (writer fills data,
sets a ready flag; reader consumes, clears it) passes data deterministically
across the boundary --- the BEL real-time pattern with sub-1\,ms latency.

%=====================================================================
\section{Q\&A Bank}
%=====================================================================

\tierbanner{Tier 1 --- Screening (phone / HR-technical)}%
{Two to four sentences.}

\question{Q1.1}{What does ``real-time'' actually mean?}
That correctness depends on \emph{timing} as well as result: the system must
respond by a deadline, predictably, even in the worst case. It does not mean
``fast'' --- a slow system that always meets its deadline is real-time; a fast
one that occasionally misses isn't.

\question{Q1.2}{Hard vs soft real-time?}
In hard real-time a missed deadline is a system failure (a flight-control loop,
an airbag) --- it must \emph{never} miss. In soft real-time an occasional miss
degrades quality but is tolerable (video, UI). The distinction drives how
conservatively you must analyse worst-case timing.

\question{Q1.3}{What are the three basic task states?}
Running (currently on the CPU), Ready (runnable but waiting for the CPU because
a higher-priority task is running), and Blocked (waiting on an event --- a
semaphore, queue, or delay). Tasks move between these as the scheduler and
events dictate.

\question{Q1.4}{What is a context switch?}
Saving the current task's CPU registers (to its stack/TCB) and restoring the
next task's, so execution resumes where that task left off. It's how one CPU
runs many tasks; on Cortex-M it's done in the PendSV exception.

\question{Q1.5}{Mutex vs semaphore in one line?}
A mutex enforces mutual exclusion and has an owner (only the locker unlocks it),
protecting a shared resource; a semaphore is an ownerless counter for signalling
events or counting resources.

\question{Q1.6}{What is a queue used for in an RTOS?}
To pass data between tasks (or from an ISR to a task) safely, with built-in
buffering and synchronisation --- the producer enqueues, the consumer blocks
until data is available. It avoids shared mutable state and its races.

\question{Q1.7}{What is preemptive scheduling?}
The scheduler always runs the highest-priority ready task, immediately
interrupting (preempting) a lower-priority one when a higher task becomes ready.
This makes worst-case response time analysable, which is why RTOSes use it.

\question{Q1.8}{Bare-metal superloop vs RTOS --- when each?}
A superloop (\code{while(1)} polling + interrupts) is simplest and fine for
small, few-task systems with loose timing. An RTOS earns its overhead when you
have multiple independent activities with different rates/priorities and need
analysable preemptive timing --- it trades RAM/complexity for structure and
determinism.

\question{Q1.9}{What is priority inversion (one line)?}
When a high-priority task is blocked waiting on a resource held by a
low-priority task, and a medium task preempts the low task --- so the
high-priority task is effectively delayed by an unrelated medium task.

\question{Q1.10}{What is rate-monotonic scheduling?}
A fixed-priority policy for periodic tasks where the shorter the period, the
higher the priority. It's optimal among fixed-priority schemes and comes with a
utilisation bound that guarantees all deadlines are met if utilisation stays
under it.

\tierbanner{Tier 2 --- Core technical (panel round)}%
{A short paragraph.}

\question{Q2.1}{Explain priority inversion fully and how priority inheritance
fixes it.}
A high task H blocks on a mutex held by a low task L. A medium task M (needing
nothing from the mutex) becomes ready and, being higher priority than L,
preempts L --- so L can't run to release the mutex, and H stays blocked behind
M indefinitely. \textbf{Priority inheritance} fixes it: while H waits on L's
mutex, the RTOS temporarily raises L's priority to H's, so L runs (M can no
longer preempt it), releases the mutex, and H proceeds; L then drops back. This
is the famous Mars Pathfinder bug --- a watchdog kept resetting because a
high-priority bus task was inverted behind a low one. The alternative,
\emph{priority ceiling}, raises the locker to a precomputed ceiling on acquire.

\question{Q2.2}{What are the four conditions for deadlock, and how do you
prevent it?}
Mutual exclusion (resources held exclusively), hold-and-wait (a task holds one
resource while waiting for another), no preemption (resources aren't forcibly
taken), and circular wait (a cycle of tasks each waiting on the next's
resource). All four must hold simultaneously, so breaking \emph{any one}
prevents deadlock --- the most practical is to eliminate circular wait by
acquiring multiple locks in a \textbf{fixed global order} everywhere. Others:
take all locks at once (no hold-and-wait), or use timeouts (so a stuck acquire
fails rather than waits forever).

\question{Q2.3}{When do you use a mutex vs a binary semaphore vs a queue?}
Mutex: protect a shared \emph{resource} that must be accessed by one task at a
time (it has ownership and supports priority inheritance, so it's the right
tool for mutual exclusion). Binary semaphore: \emph{signal} an event between
contexts (an ISR gives, a task takes) --- no ownership, used purely for
synchronisation. Queue: \emph{transfer data} with buffering, so producer and
consumer never touch shared memory directly. A common bug is using a semaphore
where a mutex is needed --- you lose priority inheritance and reintroduce
inversion.

\question{Q2.4}{Why must you use the ``FromISR'' API variants in an ISR?}
Because the normal RTOS calls may block (and an ISR must never block) and they
manipulate scheduler state assuming task context. The \code{...FromISR}
variants are non-blocking and signal the scheduler that a context switch may be
needed \emph{after} the ISR returns (via a ``higher-priority task woken'' flag
$\rightarrow$ a yield). Calling a blocking task API from an ISR corrupts the
scheduler or hangs --- a serious, common bug.

\question{Q2.5}{Apply the rate-monotonic utilisation test to three tasks.}
For tasks with (compute, period) (1,4), (2,8), (2,16), utilisation
$U = 1/4 + 2/8 + 2/16 = 0.625$ (62.5\%). The Liu--Layland bound for $n=3$ is
$3(2^{1/3}-1) \approx 0.78$. Since $0.625 \le 0.78$, the set is schedulable
under rate-monotonic priorities (shortest period highest). Note the bound is
\emph{sufficient but not necessary} --- sets above it may still be schedulable
(check with exact response-time analysis), but under it you're guaranteed safe.

\question{Q2.6}{What is the difference between a counting and a binary
semaphore, with a use for each?}
A binary semaphore holds 0 or 1 --- used to signal a single event (data ready,
interrupt occurred). A counting semaphore holds 0..N --- used to manage $N$
interchangeable resources (e.g.\ a pool of DMA buffers): each \code{take}
decrements, each \code{give} increments, and a task blocks when the count is 0.
The count \emph{is} the number of available resources, so it naturally limits
concurrency to $N$.

\question{Q2.7}{How does a queue both synchronise and transfer data?}
A task that reads an empty queue \emph{blocks} until a producer writes (the
synchronisation), and the write \emph{copies} the data into the queue's buffer
(the transfer), so the consumer gets a private copy with no shared-memory race.
A bounded queue also applies back-pressure: a producer writing to a full queue
blocks (or fails), preventing overrun. This is why RTOS designs prefer
``communicate by passing messages'' over ``share memory and lock.''

\question{Q2.8}{Describe a DPRAM flag handshake between two processors.}
The producer writes its data into the shared dual-port RAM, then sets a
\textbf{ready flag}; the consumer polls (or is interrupted on) that flag, reads
the data, and \textbf{clears} the flag to signal ``consumed.'' The flag is the
single synchronisation point, and the ordering (data \emph{then} flag on write;
flag \emph{then} data on read) plus a memory barrier guarantees the consumer
never reads half-written data. With a fixed polling period this gives bounded,
deterministic latency --- the BEL sub-1\,ms DPRAM architecture on the resume.

\question{Q2.9}{What makes an RTOS deterministic, and what can break it?}
Determinism comes from preemptive fixed-priority scheduling (the highest task
always runs within a bounded time), $O(1)$ or bounded-time scheduler operations,
and bounded interrupt latency. It's broken by: unbounded priority inversion
(fixed by inheritance), long critical sections / interrupt disables (raise worst
-case latency), dynamic memory allocation in the hot path (variable time ---
Chapter 1), and unbounded loops or blocking in high-priority tasks. Analysing
worst-case timing means accounting for all of these.

\tierbanner{Tier 3 --- Trap \& depth (principal-engineer level)}%
{A paragraph plus the trap.}

\question{Q3.1}{Your high-priority control task occasionally misses its deadline
even though average CPU load is low. Where do you look?}
Low \emph{average} load is irrelevant to a real-time miss --- you hunt the
\emph{worst case}. Suspects: (1) \textbf{priority inversion} --- the control
task blocks on a mutex (maybe hidden inside a driver or \code{printf}) held by a
low task that a medium task preempts; (2) a \textbf{long critical section or
interrupt-disable} elsewhere that delays the control task's preemption; (3) a
lower-priority ISR running long (the control task is a \emph{task}, and any ISR
outranks all tasks); (4) \textbf{blocking calls} or unbounded loops in the
path. Method: enable the RTOS's runtime stats / trace, measure worst-case
latency with a GPIO toggle on a scope, and check for mutexes without
inheritance. \trap the trap is reasoning from average utilisation; real-time is
about the worst-case confluence (the control task ready exactly when the longest
critical section and an inversion coincide). That worst-case mindset is the core
of the BEL HSI timing-defect work.

\question{Q3.2}{Is a binary semaphore a valid substitute for a mutex? When does
that choice cause a bug?}
Functionally a binary semaphore can guard a critical section, but it is
\emph{not} a mutex: it has no owner (any task can ``give'' it, even one that
didn't take it) and --- crucially --- it provides no \textbf{priority
inheritance}. So if you protect a resource with a binary semaphore and a
high-priority task waits on it, a medium task can inversion-block the low holder
indefinitely; a real mutex would have boosted the holder. \trap the trap is
``they both do mutual exclusion, so they're interchangeable.'' For resource
protection use a mutex (ownership + inheritance); reserve semaphores for
signalling. Using the wrong one reintroduces exactly the unbounded inversion
that sank Mars Pathfinder.

\question{Q3.3}{Two tasks each lock mutex A then B in one path and B then A in
another. What happens, and give two fixes?}
Classic \textbf{deadlock} via circular wait: task 1 holds A and waits for B
while task 2 holds B and waits for A --- neither can proceed. It's intermittent
because it needs the precise interleaving where each grabs its first lock before
the other grabs the second. Fix 1 (best): impose a \textbf{global lock
ordering} --- every task acquires A before B, never B before A, which makes a
cycle impossible. Fix 2: use a timeout on the second acquire and back off /
retry (release the first lock, wait, try again). \trap the trap is ``it works in
testing'' --- deadlocks are timing-dependent and may appear only in the field
under load. The disciplined design \emph{proves} deadlock-freedom by lock
ordering rather than hoping the interleaving never occurs.

\question{Q3.4}{Design the DPRAM handshake so a fast writer never corrupts what
a slower reader is consuming, and quantify the latency.}
Single-slot flag handshake: writer checks the flag is \emph{clear} before
writing (so it never overwrites unconsumed data), writes the payload, issues a
memory barrier, then sets the flag; reader sees the flag set, reads the payload,
barrier, then clears the flag. This is single-producer/single-consumer and
lock-free --- the writer owns the ``clear$\rightarrow$set'' transition, the
reader owns ``set$\rightarrow$clear.'' If the writer is faster than the reader,
it simply \emph{waits} (flag still set) and drops or buffers --- you choose the
overrun policy. Latency is bounded by the reader's poll period (or interrupt
latency): with a 1\,kHz poll, worst-case delivery is $<$1\,ms --- the resume's
figure. \trap the trap is omitting the barrier and the ``flag clear before
write'' check, which lets the reader see a torn, half-updated payload --- the
same visibility/ordering/atomicity triad from Chapters 1 and 4, now across a
processor boundary. For higher throughput you extend it to a double-buffer or a
ring in the shared RAM.

\question{Q3.5}{RMS vs EDF --- why do most safety-critical embedded systems use
fixed-priority RMS despite EDF reaching 100\% utilisation?}
EDF is dynamic (priorities change as deadlines approach) and can theoretically
schedule any set up to 100\% utilisation, but that very dynamism makes it
\emph{harder to analyse, certify, and debug}: a transient overload under EDF can
cause a \emph{domino} of missed deadlines (the whole set destabilises), and you
can't easily reason about ``which task runs when.'' Fixed-priority RMS is
predictable (priorities are static, the highest-priority task is always
identifiable), degrades gracefully under overload (only the lowest-priority
tasks miss), and has mature response-time analysis and tool/certification
support. \trap the trap is treating ``higher utilisation bound'' as ``better.''
For DO-178-style certification, \emph{analysability and predictable degradation}
beat theoretical efficiency --- which is why flight software overwhelmingly uses
fixed-priority preemptive scheduling, and why the resume's real-time work lives
in that world.

%=====================================================================
\section{Coding Gym}
%=====================================================================

\begin{partnote}
Twelve host-compilable problems (verified under
\code{gcc -Wall -Wextra -std=c11}) modelling the scheduler and synchronisation
mechanics: RMS/EDF selection, schedulability, mutex/semaphore/queue, priority
inheritance, deadlock detection, and the DPRAM mailbox handshake from the BEL
work.
\end{partnote}

%---------------------------------------------------------------
\begin{gymproblem}{Rate-monotonic priority comparison}

\textbf{Problem.} Decide which of two periodic tasks gets higher priority under
rate-monotonic scheduling.

\textbf{Hints.}
\begin{itemize}
\item Shorter period $\rightarrow$ higher priority.
\end{itemize}

\attempt

\textbf{Solution.}
\begin{cgym}
#include <stdint.h>
#include <stdio.h>

static int rms_higher(uint32_t period_a, uint32_t period_b) {
    return period_a < period_b;     /* shorter period wins */
}

int main(void) {
    printf("a(4ms) higher than b(8ms)? %d\n", rms_higher(4, 8));   /* 1 */
    return 0;
}
\end{cgym}

Rate-monotonic assigns priority purely by period: the 4\,ms task outranks the
8\,ms task, regardless of how much CPU each uses. This is provably the optimal
\emph{fixed}-priority assignment for periodic tasks, and it's why you set RTOS
task priorities by rate, not by ``importance'' --- a frequent beginner mistake
that breaks the schedulability guarantee.

\followups
\begin{enumerate}
\item \emph{``Two tasks with equal periods?''} $\rightarrow$ Tie-break
  arbitrarily (or by deadline); they can round-robin.
\item \emph{``A critical but slow task --- give it top priority?''}
  $\rightarrow$ No --- that can break RMS analysis; use period, and bound its
  execution.
\item \emph{``Deadline $\ne$ period?''} $\rightarrow$ Use deadline-monotonic
  (shorter deadline wins) --- the generalisation.
\end{enumerate}
\end{gymproblem}

%---------------------------------------------------------------
\begin{gymproblem}{Rate-monotonic schedulability (Liu--Layland)}

\textbf{Problem.} Test whether a set of periodic tasks is RMS-schedulable by the
utilisation bound.

\textbf{Hints.}
\begin{itemize}
\item $U=\sum c_i/t_i$; bound $= n(2^{1/n}-1)$.
\end{itemize}

\attempt

\textbf{Solution.}
\begin{cgym}
#include <stdint.h>
#include <stdio.h>
#include <math.h>

static int rms_schedulable(const uint32_t *c, const uint32_t *t, int n) {
    double U = 0;
    for (int i = 0; i < n; i++) U += (double)c[i] / (double)t[i];
    double bound = n * (pow(2.0, 1.0 / n) - 1.0);
    return U <= bound + 1e-9;
}

int main(void) {
    uint32_t C[3] = {1, 2, 2}, T[3] = {4, 8, 16};
    printf("schedulable=%d\n", rms_schedulable(C, T, 3));   /* 1 (U=0.625 <= 0.78) */
    return 0;
}
\end{cgym}

$U = 0.25 + 0.25 + 0.125 = 0.625$ against the $n=3$ bound $\approx 0.78$, so the
set is guaranteed schedulable. The bound is \emph{sufficient but not necessary}:
sets above it may still meet all deadlines (verify with exact response-time
analysis), but under it you need no further proof. This one calculation tells you
at design time whether a proposed task set even fits.

\followups
\begin{enumerate}
\item \emph{``Above the bound --- now what?''} $\rightarrow$ Exact
  response-time analysis; the bound is conservative.
\item \emph{``Bound as $n\to\infty$?''} $\rightarrow$ $\ln 2 \approx 0.69$ ---
  the famous 69\% rule.
\item \emph{``Where does context-switch cost go?''} $\rightarrow$ Add it into
  each task's $c_i$ for a realistic bound.
\end{enumerate}
\end{gymproblem}

%---------------------------------------------------------------
\begin{gymproblem}{Mutex with ownership}

\textbf{Problem.} Model a mutex where only the locking task may unlock.

\textbf{Hints.}
\begin{itemize}
\item Track \code{locked} and \code{owner}; reject unlock by a non-owner.
\end{itemize}

\attempt

\textbf{Solution.}
\begin{cgym}
#include <stdio.h>

typedef struct { int locked, owner; } mutex_t;

static int mutex_lock(mutex_t *m, int task) {
    if (m->locked) return 0;            /* already held -> caller must block/retry */
    m->locked = 1; m->owner = task;
    return 1;
}
static int mutex_unlock(mutex_t *m, int task) {
    if (m->locked && m->owner == task) { m->locked = 0; m->owner = -1; return 1; }
    return 0;                           /* only the owner may unlock */
}

int main(void) {
    mutex_t m = {0, -1};
    printf("lock1=%d lock2=%d unlock2=%d unlock1=%d\n",
           mutex_lock(&m, 1), mutex_lock(&m, 2),
           mutex_unlock(&m, 2), mutex_unlock(&m, 1));   /* 1 0 0 1 */
    return 0;
}
\end{cgym}

The defining mutex property is \textbf{ownership}: task 2 cannot unlock what
task 1 locked. That ownership is what lets a real RTOS implement priority
inheritance (it knows \emph{who} holds the lock to boost). A semaphore has no
owner, which is exactly why it can't provide inheritance (Q3.2). In a real RTOS
a failed lock \emph{blocks} the task rather than returning 0.

\followups
\begin{enumerate}
\item \emph{``Recursive mutex?''} $\rightarrow$ Track a count so the owner can
  re-lock and must unlock the same number of times.
\item \emph{``Why does ownership enable inheritance?''} $\rightarrow$ The RTOS
  must know which task to boost --- the holder.
\item \emph{``What should \code{lock} do on contention in a real RTOS?''}
  $\rightarrow$ Block the caller and add it to the mutex's wait list.
\end{enumerate}
\end{gymproblem}

%---------------------------------------------------------------
\begin{gymproblem}{Priority inheritance boost}

\textbf{Problem.} Compute the holder's effective priority when a higher-priority
task waits on its mutex (lower number = higher priority).

\textbf{Hints.}
\begin{itemize}
\item Boost the holder to the highest-priority (lowest-number) waiter.
\end{itemize}

\attempt

\textbf{Solution.}
\begin{cgym}
#include <stdio.h>

/* lower number = higher priority; boost holder to the waiter if waiter is higher */
static int inherit_priority(int holder_prio, int waiter_prio) {
    return waiter_prio < holder_prio ? waiter_prio : holder_prio;
}

int main(void) {
    printf("boosted=%d\n", inherit_priority(5, 2));   /* low holder(5) boosted to 2 */
    return 0;
}
\end{cgym}

When high-priority task (prio 2) blocks on a mutex held by low-priority task
(prio 5), the holder is temporarily boosted to prio 2 so a medium task (prio 3)
can no longer preempt it --- the holder runs, releases the mutex, and the high
task proceeds; the holder then reverts to 5. This is the entire mechanism that
prevents unbounded priority inversion (Q2.1, Q3.2), and the fix NASA applied to
Mars Pathfinder.

\followups
\begin{enumerate}
\item \emph{``Multiple waiters?''} $\rightarrow$ Boost to the highest-priority
  (lowest-number) waiter among them.
\item \emph{``Nested mutexes?''} $\rightarrow$ Inheritance can chain; this is
  where priority-ceiling is simpler to bound.
\item \emph{``When does the boost end?''} $\rightarrow$ On unlock the holder
  drops back to its base (or next-highest inherited) priority.
\end{enumerate}
\end{gymproblem}

%---------------------------------------------------------------
\begin{gymproblem}{Deadlock detection by cycle finding}

\textbf{Problem.} Given a wait-for graph (each task waits for at most one other),
detect a circular wait (deadlock).

\textbf{Hints.}
\begin{itemize}
\item One out-edge per node $\rightarrow$ a functional graph; use Floyd's
  tortoise/hare to find a cycle.
\end{itemize}

\attempt

\textbf{Solution.}
\begin{cgym}
#include <stdio.h>

/* wait_for[i] = task i waits for the resource held by task wait_for[i], or -1 */
static int reaches_cycle(const int *wf, int start) {
    int slow = start, fast = start;
    for (;;) {
        if (fast == -1) return 0;
        fast = wf[fast];
        if (fast == -1) return 0;
        fast = wf[fast];
        slow = wf[slow];
        if (slow == fast) return 1;     /* tortoise met hare -> cycle */
    }
}
static int has_cycle(const int *wf, int n) {
    for (int i = 0; i < n; i++) if (reaches_cycle(wf, i)) return 1;
    return 0;
}

int main(void) {
    int cyc[3] = {1, 2, 0};     /* 0->1->2->0 : deadlock */
    int ok[3]  = {1, 2, -1};    /* 0->1->2->done : no cycle */
    printf("cycle=%d nocycle=%d\n", has_cycle(cyc, 3), has_cycle(ok, 3));   /* 1 0 */
    return 0;
}
\end{cgym}

A circular wait --- task 0 waits on 1, 1 on 2, 2 back on 0 --- is the fourth
deadlock condition and the one you detect or design out. Floyd's algorithm finds
the cycle in $O(n)$ time and $O(1)$ space. In practice you \emph{prevent} rather
than detect: a global lock-acquisition order makes a cycle structurally
impossible (Q3.3), but a deadlock detector like this is useful in test/diagnostic
builds.

\followups
\begin{enumerate}
\item \emph{``How do you prevent rather than detect?''} $\rightarrow$ Fixed
  global lock ordering, or take-all-or-none, or timeouts.
\item \emph{``Multiple resources per task?''} $\rightarrow$ General
  wait-for-graph cycle detection (DFS with colours).
\item \emph{``Why is detection-after-the-fact risky in flight code?''}
  $\rightarrow$ By then you're already deadlocked; prevention is mandatory.
\end{enumerate}
\end{gymproblem}

%---------------------------------------------------------------
\begin{gymproblem}{Bounded message queue}

\textbf{Problem.} Implement a fixed-size FIFO queue for inter-task messages.

\textbf{Hints.}
\begin{itemize}
\item Circular buffer with head, tail, and count; reject on full/empty.
\end{itemize}

\attempt

\textbf{Solution.}
\begin{cgym}
#include <stdio.h>

#define QN 4
typedef struct { int buf[QN]; int head, tail, count; } queue_t;

static int q_put(queue_t *q, int v) {
    if (q->count == QN) return 0;       /* full */
    q->buf[q->tail] = v;
    q->tail = (q->tail + 1) % QN;
    q->count++;
    return 1;
}
static int q_get(queue_t *q, int *v) {
    if (q->count == 0) return 0;        /* empty */
    *v = q->buf[q->head];
    q->head = (q->head + 1) % QN;
    q->count--;
    return 1;
}

int main(void) {
    queue_t q = {{0}, 0, 0, 0};
    q_put(&q, 10); q_put(&q, 20);
    int v; q_get(&q, &v);
    printf("first=%d count=%d\n", v, q.count);   /* 10 1 */
    return 0;
}
\end{cgym}

This is the data-passing primitive: the producer \code{put}s, the consumer
\code{get}s in FIFO order, and the bounded size applies back-pressure (a full
queue rejects/blocks, preventing overrun). In an RTOS, \code{get} on an empty
queue \emph{blocks} the consumer until data arrives --- combining transfer and
synchronisation, which is why ``communicate by message passing'' beats ``share
memory and lock'' (Q2.7). It's the ring buffer of Chapter 4 with task semantics.

\followups
\begin{enumerate}
\item \emph{``Make it block instead of returning 0.''} $\rightarrow$ The RTOS
  queue suspends the task on empty/full with a timeout.
\item \emph{``Pass structs, not ints.''} $\rightarrow$ Copy by size; the queue
  owns the buffer so the sender can reuse its variable.
\item \emph{``ISR producer?''} $\rightarrow$ Use the \code{...FromISR} variant
  (Q2.4) and request a yield on wake.
\end{enumerate}
\end{gymproblem}

%---------------------------------------------------------------
\begin{gymproblem}{Round-robin: next ready task}

\textbf{Problem.} Pick the next runnable task after the current one, skipping
blocked tasks.

\textbf{Hints.}
\begin{itemize}
\item Scan circularly from \code{cur+1}; return the first ready index.
\end{itemize}

\attempt

\textbf{Solution.}
\begin{cgym}
#include <stdio.h>

static int rr_next_ready(const int *ready, int n, int cur) {
    for (int k = 1; k <= n; k++) {
        int i = (cur + k) % n;          /* wrap around */
        if (ready[i]) return i;
    }
    return -1;                          /* nothing ready -> idle */
}

int main(void) {
    int ready[4] = {0, 0, 1, 1};
    printf("next=%d\n", rr_next_ready(ready, 4, 0));   /* 2 */
    return 0;
}
\end{cgym}

Round-robin time-slices \emph{equal-priority} ready tasks for fairness: from the
current task it advances to the next ready one, wrapping around. Returning $-1$
(nothing ready) is when the scheduler runs the \textbf{idle task} (often a
low-power wait). Real RTOSes combine this with priority: round-robin only among
the tasks at the highest \emph{populated} priority level, so it's fairness
\emph{within} a priority, not across.

\followups
\begin{enumerate}
\item \emph{``Combine with priority.''} $\rightarrow$ Round-robin only the
  ready tasks at the top non-empty priority.
\item \emph{``What runs when all are blocked?''} $\rightarrow$ The idle task ---
  often \code{WFI} to sleep until an interrupt.
\item \emph{``Starvation risk?''} $\rightarrow$ A high-priority always-ready
  task starves lower ones; that's intended under strict priority.
\end{enumerate}
\end{gymproblem}

%---------------------------------------------------------------
\begin{gymproblem}{Context-switch save model}

\textbf{Problem.} Model saving a task's registers, stack pointer, and PC into
its TCB.

\textbf{Hints.}
\begin{itemize}
\item Copy the register set and the SP/PC into the control block.
\end{itemize}

\attempt

\textbf{Solution.}
\begin{cgym}
#include <stdint.h>
#include <stdio.h>

typedef struct { uint32_t r[4], sp, pc; } tcb_t;

static void ctx_save(tcb_t *t, const uint32_t *regs, uint32_t sp, uint32_t pc) {
    for (int i = 0; i < 4; i++) t->r[i] = regs[i];
    t->sp = sp;
    t->pc = pc;
}

int main(void) {
    tcb_t t;
    uint32_t regs[4] = {1, 2, 3, 4};
    ctx_save(&t, regs, 0x2000, 0x8000);
    printf("r0=%u sp=0x%X pc=0x%X\n", t.r[0], t.sp, t.pc);   /* 1 0x2000 0x8000 */
    return 0;
}
\end{cgym}

A real Cortex-M context switch is more nuanced: the hardware auto-stacks
R0--R3/R12/LR/PC/xPSR on exception entry (Chapter 3), and the PendSV handler
saves the remaining R4--R11 and swaps the stack pointer (PSP) to the next task's
--- so ``save context'' is split between hardware and the few lines of
assembly in PendSV. This model captures the \emph{concept}: each task's state
lives in its TCB/stack so the CPU can be multiplexed. The switch cost (these
saves/restores) is part of the timing budget.

\followups
\begin{enumerate}
\item \emph{``What does the hardware save vs PendSV?''} $\rightarrow$ Hardware:
  R0--R3,R12,LR,PC,xPSR; PendSV: R4--R11 and the PSP swap.
\item \emph{``Why PendSV for the switch?''} $\rightarrow$ It's the lowest
  priority, so it runs after all other ISRs --- no nested-switch races.
\item \emph{``Per-task stacks --- sizing?''} $\rightarrow$ Worst-case call depth
  + ISR frame; verify with the painting technique (Chapter 2).
\end{enumerate}
\end{gymproblem}

%---------------------------------------------------------------
\begin{gymproblem}{DPRAM mailbox handshake}

\textbf{Problem.} Implement a single-slot dual-port-RAM mailbox: a producer
writes with a ready flag, a consumer reads and clears it.

\textbf{Hints.}
\begin{itemize}
\item Producer writes only when the flag is clear; consumer reads only when set.
\end{itemize}

\attempt

\textbf{Solution.}
\begin{cgym}
#include <stdio.h>

typedef struct { volatile int flag; volatile int data; } mailbox_t;

static int mbox_write(mailbox_t *m, int v) {   /* producer side */
    if (m->flag) return 0;          /* previous value not yet consumed */
    m->data = v;                    /* write payload ... */
    m->flag = 1;                    /* ... THEN publish (barrier on real HW) */
    return 1;
}
static int mbox_read(mailbox_t *m, int *v) {   /* consumer side */
    if (!m->flag) return 0;         /* nothing ready */
    *v = m->data;                   /* read payload ... */
    m->flag = 0;                    /* ... THEN release the slot */
    return 1;
}

int main(void) {
    mailbox_t mb = {0, 0};
    mbox_write(&mb, 42);
    int got = 0;
    mbox_read(&mb, &got);
    printf("got=%d flag_after=%d\n", got, mb.flag);   /* 42 0 */
    return 0;
}
\end{cgym}

This is the BEL real-time DPRAM pattern. The flag is the single synchronisation
point between two processors sharing the dual-port RAM; the \emph{ordering} ---
data before flag on write, flag before data on read --- plus a memory barrier on
real hardware guarantees the consumer never sees a torn payload (Q3.4). With a
fixed poll period the worst-case latency is bounded by that period: a 1\,kHz poll
gives sub-1\,ms delivery, the figure on the resume. It's single-producer/
single-consumer and lock-free, like the ring buffer but across a hardware
boundary.

\followups
\begin{enumerate}
\item \emph{``Why check the flag is clear before writing?''} $\rightarrow$ So
  the producer never overwrites data the consumer hasn't taken.
\item \emph{``Higher throughput?''} $\rightarrow$ Double-buffer or a ring in the
  shared RAM (ping-pong, Chapter 4).
\item \emph{``Where's the bug if you set the flag before writing data?''}
  $\rightarrow$ The reader can grab a half-written payload --- the ordering
  trap.
\end{enumerate}
\end{gymproblem}

%---------------------------------------------------------------
\begin{gymproblem}{Counting semaphore (resource pool)}

\textbf{Problem.} Manage $N$ identical resources with a counting semaphore.

\textbf{Hints.}
\begin{itemize}
\item \code{take} decrements if available; \code{give} increments.
\end{itemize}

\attempt

\textbf{Solution.}
\begin{cgym}
#include <stdio.h>

typedef struct { int count; } csem_t;

static int  sem_take(csem_t *s) { if (s->count > 0) { s->count--; return 1; } return 0; }
static void sem_give(csem_t *s) { s->count++; }

int main(void) {
    csem_t s = {2};                         /* pool of 2 resources */
    int a = sem_take(&s), b = sem_take(&s), c = sem_take(&s);
    sem_give(&s);
    int d = sem_take(&s);
    printf("takes=%d,%d,%d give+take=%d\n", a, b, c, d);   /* 1 1 0 1 */
    return 0;
}
\end{cgym}

The count \emph{is} the number of free resources: two takes succeed, the third
fails (pool exhausted), and after a give the next take succeeds again. This
guards a pool of interchangeable resources --- DMA buffers, connection slots,
hardware channels --- limiting concurrency to $N$. In an RTOS the failing
\code{take} \emph{blocks} the task until another task \code{give}s, naturally
serialising demand on the pool. Unlike a mutex, there's no ownership: any task
may give.

\followups
\begin{enumerate}
\item \emph{``Binary semaphore from this?''} $\rightarrow$ Initialise count to 1
  (or 0 for a signal).
\item \emph{``ISR gives, task takes --- valid?''} $\rightarrow$ Yes; that's the
  classic event-signalling use (with the FromISR variant).
\item \emph{``Why is no-ownership a problem for mutual exclusion?''}
  $\rightarrow$ No priority inheritance $\rightarrow$ inversion risk (Q3.2).
\end{enumerate}
\end{gymproblem}

%---------------------------------------------------------------
\begin{gymproblem}{Earliest-deadline-first selection}

\textbf{Problem.} Among ready tasks, pick the one with the nearest absolute
deadline.

\textbf{Hints.}
\begin{itemize}
\item Scan ready tasks; track the minimum deadline.
\end{itemize}

\attempt

\textbf{Solution.}
\begin{cgym}
#include <stdint.h>
#include <stdio.h>

static int edf_next(const uint32_t *deadlines, const int *ready, int n) {
    int best = -1; uint32_t bd = 0;
    for (int i = 0; i < n; i++) {
        if (ready[i] && (best == -1 || deadlines[i] < bd)) {
            best = i; bd = deadlines[i];
        }
    }
    return best;
}

int main(void) {
    uint32_t dl[3] = {30, 10, 20};
    int ready[3]   = {1, 1, 1};
    printf("edf=%d\n", edf_next(dl, ready, 3));   /* task 1 (deadline 10) */
    return 0;
}
\end{cgym}

EDF is the dynamic-priority alternative to RMS: whoever's deadline is nearest
runs next, so priorities shift over time. It can schedule sets up to 100\%
utilisation --- more than RMS's ~69--78\% bound --- but is harder to analyse and
can cascade missed deadlines under overload (Q3.5), which is why safety-critical
fixed-priority systems usually prefer RMS. Recognising EDF (and its trade-off) is
the depth an interviewer wants.

\followups
\begin{enumerate}
\item \emph{``EDF utilisation bound?''} $\rightarrow$ 100\% on a single core for
  independent periodic tasks --- its headline advantage.
\item \emph{``Why is overload dangerous under EDF?''} $\rightarrow$ A domino of
  misses; RMS degrades more predictably (lowest priority misses first).
\item \emph{``Ties in deadline?''} $\rightarrow$ Break arbitrarily/consistently;
  doesn't affect schedulability.
\end{enumerate}
\end{gymproblem}

%---------------------------------------------------------------
\begin{gymproblem}{CPU utilisation}

\textbf{Problem.} Compute total CPU utilisation (percent) for a set of periodic
tasks.

\textbf{Hints.}
\begin{itemize}
\item $U=\sum c_i/t_i$, expressed as a percentage.
\end{itemize}

\attempt

\textbf{Solution.}
\begin{cgym}
#include <stdint.h>
#include <stdio.h>

static int cpu_util_pct(const uint32_t *c, const uint32_t *t, int n) {
    uint64_t u = 0;
    for (int i = 0; i < n; i++) u += (uint64_t)c[i] * 100u / t[i];
    return (int)u;
}

int main(void) {
    uint32_t C[3] = {1, 2, 2}, T[3] = {4, 8, 16};
    printf("util=%d%%\n", cpu_util_pct(C, T, 3));   /* 62% */
    return 0;
}
\end{cgym}

Utilisation is the fraction of CPU the periodic workload consumes ---
62\% here. It's the first feasibility number: above 100\% the set is impossible
on one core; below the RMS bound it's guaranteed schedulable; in between it
needs response-time analysis. You also leave headroom for ISRs, the context-
switch overhead, and aperiodic work --- real designs target well under 100\%.

\followups
\begin{enumerate}
\item \emph{``Above 100\% --- options?''} $\rightarrow$ Cut work, lengthen
  periods, add a core, or move work to DMA/hardware.
\item \emph{``Why integer overflow-safe with \code{uint64\_t}?''} $\rightarrow$
  \code{c*100} can exceed 32 bits for large $c$.
\item \emph{``Headroom target?''} $\rightarrow$ Often $\le$70\% to bound jitter
  and absorb aperiodic load.
\end{enumerate}
\end{gymproblem}

%=====================================================================
\section{Link to Her Resume}
%=====================================================================

\begin{resumelink}
\textbf{DPRAM real-time architecture:} ``At BEL I validated a real-time DPRAM
polling architecture with sub-1\,ms deterministic latency. That's a single-slot
flag handshake across dual-port RAM --- producer writes the payload then sets a
ready flag, consumer reads then clears it, with the ordering and a barrier
guaranteeing no torn read, and the poll period bounding worst-case latency under
1\,ms.''

\medskip
\textbf{Real-time determinism:} ``Real-time means worst-case, not average. The
timing defects I resolved are the priority-inversion / long-critical-section /
shared-data family --- I reason about the worst-case confluence and keep
critical sections and ISRs short, using a mutex (with inheritance) for resource
protection and queues for data.''

\medskip
\textbf{Scheduling judgement:} ``Safety-critical work uses fixed-priority
preemptive scheduling (rate-monotonic) because it's analysable and degrades
predictably under overload --- which matters more than EDF's higher utilisation
when you have to certify the timing.''

\medskip
\small Maps to the resume's real-time DPRAM polling (sub-1\,ms), HSI
timing-defect resolution, and RTOS (Xenomai) exposure. Keep claims to what the
resume states; the deep project scripting is Chapter 15, from
\code{inputs/INTAKE.md} only.
\end{resumelink}

%=====================================================================
\section{Self-Test Scorecard}
%=====================================================================

\begin{scorecard}
\begin{enumerate}
\item Define real-time, and hard vs soft.
\item Name the three task states and what moves a task between them.
\item Mutex vs binary semaphore vs queue --- one job each.
\item Explain priority inversion and the fix.
\item State the four deadlock conditions and the most practical prevention.
\item Apply the RMS utilisation test to (1,4),(2,8),(2,16).
\item Why is a binary semaphore a poor substitute for a mutex?
\item Describe the DPRAM flag handshake and why the write order matters.
\item Why do safety-critical systems prefer RMS over EDF despite EDF's higher
  utilisation?
\item What does the hardware save on Cortex-M exception entry vs what PendSV
  saves?
\end{enumerate}
\end{scorecard}

\begin{answerkey}
\begin{enumerate}
\item Correctness depends on meeting deadlines predictably; hard = a miss is
  failure, soft = a miss degrades quality but is tolerable.
\item Running, Ready, Blocked; the scheduler (preemption) and events
  (semaphore/queue/delay) move tasks between them.
\item Mutex = protect a resource (ownership + inheritance); binary semaphore =
  signal an event; queue = transfer data with buffering.
\item A high task blocks on a mutex held by a low task that a medium task
  preempts; priority inheritance boosts the holder to the waiter's priority.
\item Mutual exclusion, hold-and-wait, no preemption, circular wait; prevent by
  a fixed global lock-acquisition order.
\item $U=0.625 \le 0.78$ (the $n=3$ bound) $\rightarrow$ schedulable.
\item No ownership (any task can give) and no priority inheritance, so it
  reintroduces unbounded inversion for resource protection.
\item Writer fills data then sets the flag; reader sees the flag, reads, clears
  it. Data-before-flag (with a barrier) prevents the reader seeing a
  half-written payload.
\item RMS is statically analysable and degrades predictably (lowest priority
  misses first); EDF can domino under overload and is harder to certify.
\item Hardware: R0--R3, R12, LR, PC, xPSR; PendSV: R4--R11 and the PSP swap.
\end{enumerate}
\end{answerkey}

\begin{partnote}
\emph{End of Chapter 9.} Next: \textbf{Control Systems \& PID} --- open vs
closed loop, the P/I/D terms, tuning (Ziegler--Nichols), overshoot/settling/
steady-state error, and a discrete PID in C --- the control-theory foundation
behind closed-loop actuator work.
\end{partnote}
